% ============================================================
%  §4.2 Εκτίμηση μεγεθών ολίσθησης από χρονοσειρά
% ============================================================

\label{subsec:slip_estimation}

Για να τροφοδοτηθεί το μοντέλο ελαστικών σε κάθε σημείο πίστας, χρειάζονται οι γωνίες
ολίσθησης $\alpha_i$ και λόγοι ολίσθησης $\kappa_i$ κάθε τροχού. Αυτά τα μεγέθη δεν είναι
άμεσα μετρήσιμα σε κανονικές συνθήκες αγώνα (απαιτούν ειδικούς αισθητήρες), αλλά μπορούν
να εκτιμηθούν έμμεσα από δύο εύκολα διαθέσιμες ποσότητες: το προφίλ ταχύτητας $v(s)$ (από
lap sim ή τηλεμετρία) και το γεωμετρικό προφίλ πίστας $R(s)$ (ακτίνα καμπυλότητας ως
συνάρτηση θέσης).

Η διαδικασία εκτίμησης, που παρουσιάζεται σχηματικά στο \cref{fig:slip_est_flow},
βασίζεται στην αντίστροφη επίλυση του προβλήματος δυναμικής οχήματος: δεδομένης της
επιτάχυνσης, βρες ποια συνδυασμός $(\delta, \beta)$ και άρα ποιες δυνάμεις και ολισθήσεις
είναι αναγκαίες για να επιτευχθεί.

\subsubsection{Κινηματική προεπεξεργασία}

Τα ακατέργαστα δεδομένα εισόδου (προφίλ ταχύτητας CSV και curvilinear track XLSX)
επαναδειγματολη­πτούνται σε ομοιόμορφο πλέγμα 500 σημείων. Οι επιταχύνσεις υπολογίζονται
κινηματικά:

\begin{align}
    a_x(s) &= v(s) \cdot \frac{dv}{ds}  \label{eq:ax_chain2} \\
    a_y(s) &= \frac{v(s)^2}{R(s)}       \label{eq:ay_centripetal2}
\end{align}

Η αριθμητική παράγωγος $dv/ds$ είναι επιρρεπής σε θόρυβο, ιδιαίτερα σε σημεία απότομης
μεταβολής ταχύτητας. Για αυτό εφαρμόζεται φίλτρο Savitzky-Golay, που εξομαλύνει τη
χρονοσειρά διατηρώντας τη μορφή της καμπύλης (σε αντίθεση με moving average που εισάγει
phase lag).

% ----- TikZ: high-level slip estimation pipeline -----
\begin{figure}[htbp]
\centering
\begin{tikzpicture}[
    node distance=0.55cm and 1.0cm,
    block/.style   ={rectangle, draw=black, fill=white,      text width=7cm,
                     align=center, rounded corners=4pt, minimum height=0.9cm, font=\small},
    io/.style      ={rectangle, draw=black, fill=black!5,    text width=5.2cm,
                     align=center, rounded corners=2pt, minimum height=0.9cm, font=\small},
    outbox/.style  ={rectangle, draw=black, fill=black!12,   text width=5.2cm,
                     align=center, rounded corners=2pt, minimum height=0.9cm, font=\small\bfseries},
    decision/.style={diamond,   draw=black, fill=white,      text width=2.4cm,
                     align=center, font=\scriptsize, inner sep=1pt},
    arr/.style     ={->, thick, black}
]

  \node[io]                       (inp)    {Είσοδοι: γεωμετρία πίστας, χρονοσειρά ταχύτητας};
  \node[block, below=of inp]   (kinem)  {Υπολογισμός επιταχύνσεων οχήματος $\alpha_X, \alpha_Y$};
  \node[block, below=of kinem]    (startOfIteration)     {Αρχικοποίηση πλευρικής ολίσθησης\\ και γωνίας τιμονιού ($\beta_0, \delta_0$)};
  \node[block, below=of startOfIteration]    (fz)     {Υπολογισμός κατακόρυφων\\ φορτίων ανά τροχό $F_{z,i}$};
  \node[block, below=of fz]    (wheelSlips)     {Υπολογισμός ταχυτήτων και\\ μεγεθών ολίσθησης τροχών $\alpha_i,\kappa_i$};
  \node[block, below=of wheelSlips]    (tyreForces)     {Υπολογισμός πρόσφυσης\\ ελαστικών $F_{u,i}, F_{v,i}$};
  \node[block, below=of tyreForces]       (solve)  {Επίλυση εξισώσεων ισορροπίας\\ (μόνιμη κατάσταση)};
  \node[decision, below=of solve] (conv)   {σύγκλιση;};
  \node[outbox, below=of conv]    (out)    {Έξοδος: ολισθήσεις και δυνάμεις ανά τροχό ($\alpha_i, \kappa_i, F_{u,i}, F_{v,i}$)};

  \draw[arr] (inp)    -- (kinem);
  \draw[arr] (kinem)  -- (startOfIteration);
  \draw[arr] (startOfIteration)     -- (fz);
  \draw[arr] (fz)  -- (wheelSlips);
  \draw[arr] (wheelSlips)     -- (tyreForces);
  \draw[arr] (tyreForces)     -- (solve);
  \draw[arr] (solve)     -- (conv);
  \draw[arr] (conv)   -- node[right, font=\scriptsize]{ναι} (out);

  % failure branch
  \node[font=\scriptsize, right=0.8cm of conv, align=center]
    (fail) {επισήμανση\\ασύγκλιτου σημείου};
  \draw[arr, dashed] (conv.east) -- (fail.west)
      node[midway, above, font=\scriptsize]{όχι};
      
        % ----- dashed box around the iteration loop -----
  \node[draw=black, dashed, rounded corners=6pt,
        inner sep=8pt, fit=(wheelSlips)(fz)(tyreForces)(solve)] 
        (loopbox) {};
        
  \node[anchor=south, font=\normalsize, xshift=-4pt, rotate=90] 
      at (loopbox.west) {Επαναληπτική διαδικασία};

\end{tikzpicture}
\caption{Διάγραμμα ροής -- εκτίμηση ολισθήσεων και δυνάμεων ανά τροχό από δεδομένα γύρου.}
\label{fig:slip_est_flow}
\end{figure}
% ----- τέλος flowchart -----


\subsubsection{Επίλυση 2-DOF}

Για κάθε σημείο $s_k$, δεδομένου $(v_k, a_{x,k}, a_{y,k})$, επιλύεται το σύστημα
\cref{eq:force_balance} για τους αγνώστους $\delta_k$ και $\beta_k$. Ο solver
(\texttt{lsqnonlin}) ελαχιστοποιεί το υπόλοιπο $\|\sum_i \mathbf{F}_i - m\mathbf{a}\|^2$
με φυσικά bounds.

Στο εσωτερικό κάθε κλήσης residual, η υπολογιστική αλυσίδα περιλαμβάνει: μετασχηματισμούς
πλαισίων $\to$ κατανομή $F_z$ $\to$ ταχύτητες τροχών $\to$ ολισθήσεις $\to$ δυνάμεις
ελαστικού $\to$ global frame. Αυτή η αλυσίδα εκτελείται αρκετές φορές ανά σημείο
(iterations του solver), οπότε η αποδοτικότητα κάθε συνάρτησης είναι κρίσιμη για τον
συνολικό χρόνο εκτέλεσης.

Σε σημεία ακραίων μανουβρών (π.χ. συνδυασμός επιτάχυνσης και στροφής στα όρια), το σύστημα
ενδέχεται να μη συγκλίνει εντός των ορίων. Τέτοια σημεία σημαίνονται (flag) και
διαδίδονται στο θερμικό integrator με τις τελευταίες συγκλίνουσες τιμές.

% TODO: Συζήτησε τον αντίκτυπο των non-converged σημείων στα θερμικά αποτελέσματα.
% Είναι αρκετά σπάνια ώστε να μη δημιουργούν πρόβλημα; Αξίζει να δείξεις
% χάρτη converged/unconverged κατά μήκος της πίστας.
